\documentclass[a4paper,12pt]{article}
\usepackage[utf8]{inputenc}
\usepackage[T1]{fontenc}
\usepackage[ngerman]{babel}
\usepackage{geometry}
\geometry{margin=2.5cm}

\begin{document}

\section*{PDEBench: An Extensive Benchmark for Scientific Machine Learning}

\textbf{Autoren:} Makoto Takamoto, Timothy Praditia, Raphael Leiteritz, Dan MacKinlay,
Francesco Alesiani, Dirk Pflüger, Mathias Niepert \\
\textbf{Jahr:} 2022 \\
\textbf{Konferenz:} NeurIPS 2022 Track on Datasets and Benchmarks

\subsection*{Zusammenfassung}

PDEBench stellt einen umfassenden und öffentlich zugänglichen Benchmark für maschinelles
Lernen im Bereich partieller Differentialgleichungen (PDEs) bereit. Motivation des Werkes
ist die Feststellung, dass trotz erheblicher Fortschritte im wissenschaftlichen maschinellen
Lernen (Scientific ML) ein standardisierter, gleichzeitig praktikabler und ausreichend
anspruchsvoller Benchmark fehlt, der eine faire und reproduzierbare Vergleichsbasis für
neue Methoden bietet.

\subsubsection*{Datensätze und abgedeckte PDEs}

Der Benchmark umfasst insgesamt 11 verschiedene PDE-Typen in ein-, zwei- und
dreidimensionalen Raumgebieten, die sowohl zeitabhängige als auch stationäre Probleme
abdecken. Zu den grundlegenderen Gleichungen zählen die 1D-Advektionsgleichung, die
Burgers-Gleichung, Diffusions-Reaktions-Gleichungen sowie die Diffusions-Sorptions-Gleichung
und der 2D-Darcy-Fluss. Als anspruchsvollere, praxisnahe Probleme werden zudem die
kompressible und inkompressible Navier-Stokes-Gleichung sowie die Flachwassergleichungen
berücksichtigt. Diese Auswahl umfasst eine breite Palette physikalischer Phänomene,
darunter Schockwellenbildung, turbulente Dynamik, sensitiv von Anfangsbedingungen
abhängige Systeme sowie verschiedene Randbedingungen (periodisch, Dirichlet, Neumann,
Cauchy). Die bereitgestellten Datensätze wurden durch hochauflösende numerische Simulationen
erzeugt und im HDF5-Format gespeichert; sie sind über das DaRUS-Datenrepositorium der
Universität Stuttgart dauerhaft unter einem DOI zugänglich.

\subsubsection*{Baseline-Modelle}

Die Autoren evaluieren drei etablierte ML-Architekturen als Baselines:
\begin{itemize}
    \item \textbf{Fourier Neural Operator (FNO):} Ein neuronaler Operator, der im
    Fourierraum arbeitet und dadurch auflösungsunabhängige Vorhersagen ermöglicht.
    FNO erzielt auf den meisten betrachteten Problemen die besten Ergebnisse, insbesondere
    bei glatten Strömungsfeldern, da er differentielle Operatoren im Spektralraum effizient
    approximiert.
    \item \textbf{U-Net:} Eine Encoder-Decoder-Architektur mit Skip-Connections, die
    ursprünglich für medizinische Bildsegmentierung entwickelt wurde und hier für
    PDE-Surrogate eingesetzt wird. U-Net eignet sich besonders für lokal geprägte
    Strukturen und diffusionsartige Aufgaben, zeigt jedoch Instabilitäten im vollständig
    autoregressiven Betrieb. Die Autoren empfehlen daher das Training mit dem sogenannten
    \emph{Pushforward Trick}, der die Langzeitstabilität deutlich verbessert.
    \item \textbf{Physics-Informed Neural Networks (PINNs):} Netze, die den PDE-Residual
    direkt in der Verlustfunktion minimieren und damit physikalische Konsistenz erzwingen.
    PINNs sind jedoch auf einzelne Anfangs- und Randbedingungen beschränkt und erfordern
    für jede neue Konfiguration ein vollständiges Neutraining.
\end{itemize}

\subsubsection*{Evaluierungsmetriken}

Ein wesentlicher Beitrag von PDEBench liegt in der Einführung erweiterter Bewertungsmetriken,
die über den klassischen RMSE hinausgehen. Neben dem normierten RMSE (nRMSE) und dem
maximalen Fehler werden physikalisch motivierte Metriken verwendet: der \texttt{cRMSE}
misst die Abweichung von erhaltenen physikalischen Größen, der \texttt{bRMSE} bewertet
die Einhaltung von Randbedingungen, und der \texttt{fRMSE} analysiert den Fehler im
Fourierraum getrennt nach niedrigen, mittleren und hohen Frequenzen. Diese Differenzierung
ist insbesondere in turbulenten oder unstetigen Strömungsregimen notwendig, da ein globaler
RMSE lokale Fehlerstrukturen verschleiert.

\subsubsection*{Zentrale Ergebnisse}

Die Auswertungen zeigen, dass keine einzelne Architektur auf allen Problemen dominiert:
FNO übertrifft U-Net bei den meisten Metriken, während U-Net bei diffusionsähnlichen
Aufgaben wie dem Darcy-Fluss konsistent bessere Ergebnisse liefert. Bei hochgradig
nichtlinearen oder inviskosen Strömungen versagen beide Methoden in erheblichem Maße,
was auf fundamentale Limitierungen der aktuellen Surrogatmodelle hinweist. Darüber hinaus
zeigt FNO bei abnehmenden Diffusionskoeffizienten stark erhöhte Fehler im hochfrequenten
Bereich, was auf das Gibbs-Phänomen zurückgeführt wird. Hinsichtlich der Recheneffizienz
ist die Inferenzzeit der ML-Modelle bis zu drei Größenordnungen geringer als die der
klassischen numerischen Löser, wobei die Trainingszeit erheblich ist.

\subsection*{Bezug zur Masterarbeit}

PDEBench ist für die vorliegende Masterarbeit aus mehreren Perspektiven relevant.
Erstens bietet es eine methodische Orientierung für den systematischen Vergleich von
U-Net und FNO als Surrogatmodelle, der auch in dieser Arbeit im Zentrum steht.
Die in PDEBench beobachtete komplementäre Stärkenverteilung beider Architekturen ---
FNO für globale Muster im Spektralraum, U-Net für lokale Randeffekte --- motiviert
die differenzierte Betrachtung im Kontext von Stokes-Strömungen um sedimentierende
Kristalle, wo sowohl fernfeldige Wechselwirkungen als auch lokale Störungen durch
Kristallränder eine wesentliche Rolle spielen.

Zweitens liefert PDEBench eine Vorlage für das Evaluierungsframework: Die Verwendung
physikalisch motivierter Metriken, insbesondere die Überprüfung der Divergenzfreiheit
des Geschwindigkeitsfeldes (analog zu \texttt{cRMSE}) und die Frequenzraumanalyse über
\texttt{fRMSE}, findet eine direkte Entsprechung in den 47 Metriken der Masterarbeit,
die ebenfalls über den reinen MAE hinausgehen und Größen wie Divergenzverletzungen,
Wirbelstärkeerhaltung und Randbedingungserfüllung einschließen.

Drittens unterstreicht die in PDEBench dokumentierte Schwäche beider Architekturen bei
der zeitlichen Extrapolation und bei stark nichtlinearen Regimen die Bedeutung des in
der Masterarbeit gewählten Ansatzes des \emph{Residual Learning}: Statt die gesamte
Strömung direkt zu lernen, wird das neuronale Netz auf die Abweichung von einer
analytischen Stokes-Lösung trainiert, was die Aufgabe erheblich vereinfacht und die
physikalische Konsistenz durch den analytischen Anteil absichert. Dieser Ansatz ist
in PDEBench nicht vertreten und stellt damit einen originären Beitrag der Masterarbeit
dar.

Schließlich bestätigt PDEBench die grundsätzliche Machbarkeit des Einsatzes von
FNO und U-Net als Surrogatmodelle für Strömungssimulationen, während die
beobachteten Schwächen bei der Generalisierung auf ungesehene Parameterregionen ---
insbesondere auf unterschiedliche Kristallanzahlen --- die zentrale Forschungsfrage
der Masterarbeit motivieren und deren Relevanz unterstreichen.

\end{document}